\section{Example: ''Real'' Beams, Part I} \label{ex6}

Up to now, we have only considered single electron beams, except in example
\ref{ex3}, where we have seen the effect of the energy-spread of the
beam. In this example, we learn about the first of two methods how \W~takes the
effects of real beams, i.e. emittance and enery-spread into account.\\

In modern low emittance storage rings the spatial and spectral
distribution of the synchrotron radiation does not vary very much for
the different electrons in the beam phase-space. Thus, the shape of
the radiation cone is more or less the same, it is only shifted
according to the position and slope of the considered electron.
Therefore, it is possible to take the beam emittance into account by
folding the radiation cone of the reference electron with
a distribution that reflects the phase-space. This distribution is
usually a 2D-Gaussian with
\begin{eqnarray*}
\Sigma_y^2 &=& \sigma_y^2 + d^2 \times \sigma_y'^2   \\
\Sigma_z^2 &=& \sigma_x^2 + d^2 \times \sigma_x'^2
\end{eqnarray*}

where $\sigma_y,\, \sigma_y', \sigma_z, \, \sigma_z'$ denote the
vertical and horizontal sizes
and divergences of the beam  and $d$ the distance between the source and the
observation plane. Enter:\\

\cboxs{\senv\\ \grau
python3 ../python/wave\_example\_6.py\\
cp ../input/wave\_example\_6.in wave.in\\
wave\\
wplot\\
}

Use the plotting GUI to produce the Fig. \ref{ex5fig4} or enter in the
python shell:

\cboxs{zone(1,2,1)\\
ndistpinh()\\
nextzone()\\
ndistpinv()
}

If you want to change parameters for varies runs, but
plot always the same, you can put these lines into a file

\gtis{waveplot\_startup.py}

The commands in this file are exceuted at the start of \wple, until the
end of the file is reached, or a line with a key-word \gtis{EOF} is
found.
Of course, this can also handle the plotting zones by using the \gtis{Options} menu of
\wple.

\begin{figure}[htb]
    \centering
    %\framebox(10,8){
          \includegraphics[width=0.6\textwidth]
	  {/home/scheer/wav/doc/pp/wave_example_6_hvcut.pdf}
    %} % end framebox
    \caption{Horizontal and vertical cuts through the
    flux-density distribution}
    \label{ex5fig4}
\end{figure}

The mean value of the histgrams agree with the position of the
observation pinhole as given in \woe. We see, that the position of the
pinhole is shifted. This has been achieved by setting\\

\cboxss{YSTART=56.1e-6\\
ZSTART=648.e-6\\
VYIN=12.87e-6\\
VZIN=74.4e-6\\
PINCEN(2)=-9999.\\
PINCEN(3)=-9999.\\
}

in \wie . The values for the start conditons of the electron correspond
to three sigmas of the beam sizes and divergences as given in the
namelist \scap{\$WFOLDN}. The values for the center of the pinhole
trigger the shift of the pinhole, as is explained in the namelist
\scap{\$PINHOLE}.

%\newpage
When we set in \wi\\

\cboxss{YSTART=0.0\\
YSTART=0.0\\
ZSTART=0.0\\
VYIN=0.0\\
VZIN=0.0\\
}

as usual, we get nearly the same histograms, except that they centered
on-axis. Hence, the ansatz of the folding procedure is verified.

A similar folding procedure is applied for the spectrum to take into
account the energy-spread of the beam (\scap{ESPREAD} in the namelist
\scap{\$WFOLDN}). This procedure makes use of the fact, that the
harmonics of a undulator depend on the square of the beam energy
$E_{beam}$. Thus, the shift of the spectrum is in first order given
by\\[-1.5cm]

\begin{eqnarray*}
\frac{\Delta E_{phot}}{E_{phot}}  &=&  2 \times \frac{\Delta E_{beam}}{E_{beam}}
\end{eqnarray*}
while the shape stays in a good approximation the same.

\underline{Proposed exercise:}\\

Try to verify the last statement by using \W.
